\documentclass[12pt]{article}
\usepackage[margin=2.2cm]{geometry}
\usepackage{fontspec}
\setmainfont{DejaVu Serif}
\usepackage{parskip}
\usepackage{xcolor}
\begin{document}
\section*{Apêndice de atualização -- ciclo C29}
\textbf{Decisão incorporada.} Na tela comparativa entre categorias, a mensagem explicativa da faixa superior foi substituída pela formulação:\\

{\large\bfseries No Gérard, a mesma solução numérica não implica a mesma estrutura matemática.}\\
{\large A construção dos diagramas torna visíveis essas distinções entre categorias.}

\textbf{Classificação do ciclo.} Refinamento visual e operacional.

\textbf{Justificativa.} A alteração não modifica a arquitetura, o modelo de domínio, a persistência, a classificação das categorias nem as regras de sincronização entre representações. Seu papel é reforçar, com maior hierarquia visual, a tese central que já orientava o modo comparativo.

\textbf{Impacto na complexidade do Gérard.} O acréscimo de complexidade é baixo e localizado. A mudança concentrou-se em dois pontos: (i) atualização das chaves de internacionalização da explicação da tela comparativa; e (ii) ajuste da composição visual do cabeçalho dessa mesma tela. Não houve criação de novos objetos de domínio, novos fluxos de interação, nem novos tipos de log. Em termos analíticos, trata-se de um refinamento de comunicação pedagógica, e não de ampliação estrutural.

\textbf{Efeito pedagógico esperado.} A interface passa a explicitar com maior força argumentativa que soluções numericamente equivalentes podem corresponder a estruturas matemáticas distintas, tornando mais evidente o motivo pelo qual o Gérard solicita a construção de diagramas diferentes para categorias diferentes.
\end{document}
